\section{Đánh giá hiện trạng và khoảng trống năng lực}

\subsection{Bằng chứng hiện trạng tại Đại học Bách khoa Hà Nội}
Để bảo đảm tính khách quan và liêm chính học thuật, hiện trạng công tác cố vấn học tập và cảnh báo học vụ tại Đại học Bách khoa Hà Nội được tổng hợp thông qua sự kết hợp giữa dữ liệu thứ cấp được công bố chính thức \citep{phungnhung2022}, văn bản quy chế đào tạo tín chỉ hiện hành \citep{hust2023}, và các giả định tham số hóa bối cảnh của nhóm nghiên cứu.

Xét trên phương diện hạ tầng kỹ thuật và dữ liệu, hệ thống thông tin sinh viên của trường hiện lưu trữ hơn 35.000 hồ sơ sinh viên và cập nhật điểm học phần theo từng học kỳ. Hệ thống quản trị học tập ghi nhận trung bình hơn 200.000 lượt tương tác mỗi tuần, nhưng toàn bộ nguồn dữ liệu này mới chỉ được lưu trữ dưới dạng nhật ký thô mà chưa được trích xuất để phân tích tương quan với kết quả học tập. Đáng lưu ý là giữa hai hệ thống này chưa có giao diện lập trình ứng dụng hay đường ống dữ liệu tự động kết nối theo thời gian thực; việc chuyển giao điểm số từ hệ thống quản trị học tập sang hệ thống quản lý đào tạo vẫn thực hiện thủ công thông qua việc xuất và nhập các tệp bảng tính Excel.

Đối với khía cạnh công tác cố vấn học tập, mỗi giảng viên đảm nhiệm vai trò cố vấn hiện phụ trách trung bình từ 60 đến 100 sinh viên. Theo quy chế đào tạo hiện hành, cố vấn học tập chủ yếu thực hiện nhiệm vụ họp lớp từ một đến hai lần mỗi kỳ nhằm phổ biến quy chế đăng ký học tập. Khi nhận danh sách sinh viên bị cảnh báo học vụ sau khi đã có kết quả thi cuối kỳ, cố vấn học tập chỉ có thể lập biên bản nhắc nhở mà không có bất kỳ công cụ sư phạm nào để hỗ trợ phục hồi điểm số cho học kỳ đã khép lại.

Về tỷ lệ sa sút học vụ, bên cạnh con số 700 đến 800 sinh viên bị buộc thôi học mỗi năm đã được công bố chính thức \citep{phungnhung2022}, mô hình nghiên cứu tình huống giả định tỷ lệ sinh viên chớm nhận cảnh báo học vụ mức một và mức hai chiếm khoảng 5\% đến 8\% tổng số sinh viên chính quy toàn trường, trong đó hơn 65\% số trường hợp rơi vào sinh viên năm thứ nhất và năm thứ hai do nợ các học phần toán và vật lý đại cương. Thời điểm phát hiện sinh viên có nguy cơ thường diễn ra ở tuần thứ 17 đến tuần thứ 19, tức khoảng ba đến bốn tuần sau khi kết thúc đợt thi cuối kỳ. Tại thời điểm này, sinh viên đã tích lũy điểm không đạt và điểm trung bình tích lũy đã tụt xuống dưới ngưỡng xử lý học vụ.

\subsection{Bảng đánh giá trưởng thành kép}
Để đánh giá toàn diện năng lực của nhà trường, nghiên cứu áp dụng phương pháp Đánh giá trưởng thành kép, trình bày song song hai khung đánh giá độc lập gồm Smart University Maturity Model (Mô hình Trưởng thành Đại học Thông minh - SMM) do \citet{uskov2018} đề xuất gồm năm cấp độ và Khung trưởng thành số giáo dục đại học của \citet{amiri2026} gồm tám chiều năng lực được đánh giá theo bốn cấp độ.

Quá trình đánh giá tuân thủ nguyên tắc không quy đổi cơ học giữa hai thang đo do hai khung có đối tượng và logic đánh giá khác nhau. Khung SMM phản ánh mức độ thể chế hóa và năng lực thông minh của hệ thống công nghệ, trong khi khung của Amiri và cộng sự phản ánh năng lực vận hành tổ chức số toàn diện.

\vspace{0.3cm}
\noindent\textbf{Bảng 4.1. Đánh giá trưởng thành kép cho Dịch vụ Cố vấn học tập và Hỗ trợ sớm}
\vspace{0.1cm}

\begin{table}[htbp]
\centering
\fontsize{8.5}{10.5}\selectfont
\begin{tabularx}{\textwidth}{|L{2.0cm}|L{2.8cm}|C{1.1cm}|C{1.1cm}|C{1.4cm}|C{1.1cm}|Y|}
\hline
\textbf{Khung đánh giá} & \textbf{Chiều năng lực / Đặc tính} & \textbf{AS-IS} & \textbf{TO-BE} & \textbf{Tin cậy} & \textbf{Ưu tiên} & \textbf{Minh chứng thực tế hỗ trợ đánh giá} \\
\hline
\multirow{6}{2.0cm}{\textbf{SMM (Uskov et al., 2018)}\\\textit{(Thang 1--5)}} 
& \textbf{01. Cảm nhận (Sensing)} & \textbf{2} & \textbf{4} & Cao & Cao & Thu thập nhật ký truy cập LMS và điểm danh, chưa đồng bộ thời gian thực. \\
\cline{2-7}
& \textbf{02. Suy diễn (Inferring)} & \textbf{1} & \textbf{4} & TB & Cao & Chưa có công cụ suy diễn; cố vấn tự suy đoán nguyên nhân khi sinh viên giải trình. \\
\cline{2-7}
& \textbf{03. Tự học (Learning)} & \textbf{1} & \textbf{3} & Cao & Cao & Chưa có mô hình học từ dữ liệu các khóa trước; phân tích hoàn toàn thủ công. \\
\cline{2-7}
& \textbf{04. Dự đoán (Anticipating)} & \textbf{1} & \textbf{4} & Cao & Cao & Hoàn toàn không có dự đoán; cảnh báo mang tính hồi cứu sau khi kỳ thi kết thúc. \\
\cline{2-7}
& \textbf{05. Thích ứng (Adapting)} & \textbf{2} & \textbf{3} & TB & TB & Một số giảng viên tự điều chỉnh đề cương phụ đạo, chưa có cơ chế cấp trường. \\
\cline{2-7}
& \textbf{06. Tự tổ chức (Self-org.)} & \textbf{1} & \textbf{3} & Thấp & TB & Phân bổ cố vấn theo niên chế cố định; chưa tự động phân bổ theo mức độ rủi ro. \\
\hline
\multirow{8}{2.0cm}{\textbf{Amiri et al. (2026)}\\\textit{(Thang 1--4)}} 
& \textbf{01. Quản trị chiến lược} & \textbf{2} & \textbf{3} & Cao & TB & Chiến lược chuyển đổi số đã ban hành, chưa có chính sách phân tích học tập. \\
\cline{2-7}
& \textbf{02. Văn hóa số} & \textbf{2} & \textbf{3} & TB & Cao & Sẵn sàng dùng phần mềm, còn e ngại dữ liệu giám sát và định kiến dán nhãn. \\
\cline{2-7}
& \textbf{03. Năng lực số} & \textbf{2} & \textbf{3} & Cao & Cao & Thành thạo phần mềm dạy học, thiếu kỹ năng khai thác bảng điều khiển phân tích. \\
\cline{2-7}
& \textbf{04. Hạ tầng số} & \textbf{3} & \textbf{3} & Rất cao & TB & Hạ tầng mạng, máy chủ, tài khoản Microsoft Office 365 và Moodle phủ sóng toàn trường. \\
\cline{2-7}
& \textbf{05. Dữ liệu \& Phân tích} & \textbf{2} & \textbf{3} & Rất cao & Rất cao & Dữ liệu lớn nhưng nằm ở các ốc đảo biệt lập; thiếu kiến trúc dữ liệu tập trung. \\
\cline{2-7}
& \textbf{06. Dạy và học} & \textbf{2} & \textbf{3} & Cao & Cao & Đào tạo kết hợp phổ biến; các hoạt động đánh giá quá trình chưa chuẩn hóa nhập điểm. \\
\cline{2-7}
& \textbf{07. Vận hành dịch vụ} & \textbf{2} & \textbf{3} & TB & Cao & Dịch vụ sinh viên một cửa vận hành, chưa tích hợp dòng can thiệp của cố vấn. \\
\cline{2-7}
& \textbf{08. Nghiên cứu \& Đổi mới} & \textbf{2} & \textbf{3} & Thấp & Thấp & Có các nhóm nghiên cứu trí tuệ nhân tạo mạnh nhưng chưa ứng dụng vào quản trị đào tạo. \\
\hline
\end{tabularx}
\begin{flushleft}
\textit{\fontsize{8}{10}\selectfont Ghi chú: Quy ước cấp độ SMM (1-Khởi phát, 2-Có thể lặp lại, 3-Được chuẩn hóa, 4-Được quản trị định lượng, 5-Tối ưu hóa liên tục). Quy ước cấp độ Amiri et al. (1-Mới nổi, 2-Thiết lập, 3-Tăng cường, 4-Trưởng thành).}
\end{flushleft}
\end{table}

\subsection{Hai Nhận định Ưu tiên Cốt lõi và Điều kiện Nền tảng}
Từ kết quả đối chiếu trên bảng đánh giá trưởng thành kép và phân tích ma trận khoảng cách, nghiên cứu rút ra hai nhận định ưu tiên mang tính quyết định đối với sự thành công của dịch vụ.

Nhận định ưu tiên thứ nhất chỉ ra rằng điểm nghẽn cổ chai lớn nhất hiện nay nằm ở năng lực suy diễn, dự đoán và sự phân mảnh của dữ liệu. Mặc dù nhà trường đã đạt mức Tăng cường về hạ tầng số với hệ thống máy chủ và mạng diện rộng sẵn sàng, nhưng năng lực suy diễn và dự đoán của hệ sinh thái đào tạo vẫn chỉ dừng ở mức Khởi phát. Nguyên nhân gốc rễ là nguồn dữ liệu hành vi trên hệ thống quản trị học tập và dữ liệu học vụ trên hệ thống thông tin sinh viên bị giam cầm trong các cơ sở dữ liệu riêng rẽ. Việc tiếp tục đầu tư các mô hình trí tuệ nhân tạo phức tạp sẽ không mang lại giá trị nếu chưa giải quyết được việc đồng bộ luồng dữ liệu thời gian thực. Do đó, ưu tiên kỹ thuật số một trong giai đoạn khởi đầu là xây dựng đường ống tích hợp dữ liệu và động cơ phân tích rủi ro học tập.

Nhận định ưu tiên thứ hai khẳng định sự lệch pha nghiêm trọng giữa năng lực công nghệ và văn hóa quy trình can thiệp sư phạm. Sự chênh lệch giữa công nghệ sẵn có và khả năng hành động thực tế là một rủi ro lớn. Nếu hệ thống công nghệ đưa ra cảnh báo sớm ở tuần thứ ba nhưng nhà trường không thiết lập một quy trình bàn giao trách nhiệm rõ ràng và cố vấn học tập không được bồi dưỡng kỹ năng tư vấn bối cảnh, các tín hiệu cảnh báo sẽ nhanh chóng bị bỏ qua do hiện tượng trơ lì cảnh báo. Vì vậy, ưu tiên tổ chức hàng đầu cần triển khai song song với kỹ thuật là thể chế hóa sổ tay quy trình hỗ trợ sớm và ban hành chính sách công nhận giờ làm việc cho cố vấn học tập khi thực hiện tư vấn cá nhân hóa.

Để bảo đảm việc mở rộng dịch vụ diễn ra bền vững, nhà trường bắt buộc phải thỏa mãn hai điều kiện nền tảng. Thứ nhất, trường cần hoàn thiện chính sách quản trị và bảo vệ dữ liệu cá nhân thông qua việc ban hành quy chế sử dụng dữ liệu học tập có trách nhiệm tuân thủ Nghị định 13/2023/NĐ-CP, bảo đảm sinh viên có quyền kiểm soát thông tin cá nhân. Thứ hai, trường cần chuẩn hóa quy trình đánh giá quá trình bằng việc yêu cầu giảng viên bộ môn cam kết cập nhật điểm danh và tối thiểu hai đầu điểm đánh giá quá trình hoặc tiến độ nộp bài trước tuần thứ năm của học kỳ nhằm cung cấp dữ liệu đầu vào chuẩn xác cho hệ thống phân tích.
